\section*{Hoofdstuk \ref{ch:g7:earthquakes} --- Aardbevingen}

\begin{solution}{exo:g7:earthquakes:1}
(a) \omterm{def:g7:earthquakes:focus}{haard}; (b) \omterm{def:g7:earthquakes:intensity}{intensiteit}; (c) \omterm{def:g7:earthquakes:fault}{breuk}; (d) \omterm{def:g7:earthquakes:focus}{epicentrum}; (e) \omterm{def:g7:earthquakes:magnitude}{magnitude};
(f) \omterm{def:g7:earthquakes:seismometer}{seismometer}.
\end{solution}

\begin{solution}{exo:g7:earthquakes:2}
(a) In het dorp zelf: het \omterm{def:g7:earthquakes:focus}{epicentrum} is het punt aan het oppervlak recht
boven de \omterm{def:g7:earthquakes:focus}{haard}. (b) \qty{30}{km}, recht naar beneden. (c) Meestal lager:
het schudden wordt zwakker met de afstand. (Het hangt ook af van de
ondergrond: op zachte modder zou het tweede dorp toch hevig kunnen
schudden.)
\end{solution}

\begin{solution}{exo:g7:earthquakes:3}
(a) \omterm{def:g7:earthquakes:magnitude}{Magnitude} (een getal met één decimaal). (b) \omterm{def:g7:earthquakes:intensity}{Intensiteit} (wat er op één
plek werd gezien). (c) \omterm{def:g7:earthquakes:intensity}{Intensiteit} (Romeinse cijfers, verschillend op twee
plekken). (d) \omterm{def:g7:earthquakes:magnitude}{Magnitude}: de vrijgekomen energie, één stap in \omterm{def:g7:earthquakes:magnitude}{magnitude}
meer.
\end{solution}

\begin{solution}{exo:g7:earthquakes:4}
(a) Niet waar: wetenschappers kunnen \omterm{def:g4:volcanoes-earthquakes-first-look:earthquake}{aardbevingen} niet voorspellen; ze
weten alleen welke \omterm{def:g7:earthquakes:fault}{breuken} ooit opnieuw zullen verschuiven. (b) Niet waar:
op de open \omterm{def:g2:water-on-earth:ocean}{oceaan} is een \omterm{def:g7:earthquakes:tsunami}{tsunami} zelden hoger dan een meter; hij groeit bij
de \omterm{def:g1:landscapes-around-us:water}{kust}. (c) Waar. (d) Waar: 92~\% van de \omterm{def:g4:volcanoes-earthquakes-first-look:earthquake}{aardbevingen} met \omterm{def:g7:earthquakes:magnitude}{magnitude}~6 of
meer vindt plaats op minder dan \qty{300}{km} van een \omterm{def:g7:plate-tectonics:plate-boundary}{plaatgrens}.
\end{solution}

\begin{solution}{exo:g7:earthquakes:5}
(a) Ongeveer 127 per jaar met \omterm{def:g7:earthquakes:magnitude}{magnitude} 6 tot 6.9; ongeveer 13 met
\omterm{def:g7:earthquakes:magnitude}{magnitude} 7 tot 7.9. (b) $126.7 \div 13.12 \approx 9.7$: ongeveer 10~keer
minder. (c) $13.12 + 0.76 + 0.04 = 13.92$: ongeveer \textbf{14} per jaar.
(d) $0.04 \times 50 = 2$: twee in 50~jaar (Sumatra 2004 en Japan 2011).
\end{solution}

\begin{solution}{exo:g7:earthquakes:6}
(a) Eén stap: 32~keer. (b) Twee stappen: $32 \times 32 = 1024$, ongeveer
1000~keer. (c) Elke stap maakt de golven 10~keer groter: $10 \times 10 =
100$~keer groter.
\end{solution}

\begin{solution}{exo:g7:earthquakes:7}
(a) $8 \times 5 = \qty{40}{km}$. (b) $8 \times 25 = \qty{200}{km}$. (c)
$120 \div 8 = \qty{15}{s}$. (d) Eén station geeft alleen een afstand, geen
richting: alle punten op die afstand vormen een cirkel rond het station,
en het \omterm{def:g7:earthquakes:focus}{epicentrum} kan overal op die cirkel liggen.
\end{solution}

\begin{solution}{exo:g7:earthquakes:8}
(a) V. (b) II. (c) VIII. (d) De \omterm{def:g7:earthquakes:intensity}{intensiteit} verschilt van plek tot plek:
verslag (c) komt van een plek dicht bij het \omterm{def:g7:earthquakes:focus}{epicentrum}, of gebouwd op
zachte grond; verslag (a) van een plek verder weg, of op \omterm{def:g3:stones-pebbles-sand:rock}{vast gesteente}. De
\omterm{def:g4:volcanoes-earthquakes-first-look:earthquake}{aardbeving} heeft toch maar één \omterm{def:g7:earthquakes:magnitude}{magnitude}.
\end{solution}

\begin{solution}{exo:g7:earthquakes:9}
(a) $7031 - 5561 = 1470$. (b) $1470 \div 7031 \approx 0.209$: ongeveer
\textbf{21~\%}. (c) De zakkende \omterm{def:g7:plate-tectonics:plate}{plaat} duikt diep de \omterm{def:g6:inside-the-earth:layers}{mantel} in. Ze is nog
stijf genoeg om te breken, dus kunnen er \omterm{def:g4:volcanoes-earthquakes-first-look:earthquake}{aardbevingen} in ontstaan,
honderden kilometers diep.
\end{solution}

\begin{solution}{exo:g7:earthquakes:10}
(a) $2500 \div 800 = 3.125$~uur: 3~uur en $0.125 \times 60 = 7.5$
minuten, ongeveer 3~uur 8~minuten. (b) 3~u~8~min $-$ 2~u~17~min $=$
51~minuten later. (c) De golf deed er uren over om aan te komen, terwijl
een grote \omterm{def:g4:volcanoes-earthquakes-first-look:earthquake}{aardbeving} onder de \omterm{def:g2:water-on-earth:ocean}{zee} binnen enkele minuten wordt
gelokaliseerd: een waarschuwing had de mensen naar hoger gelegen grond
kunnen sturen, lang voordat de golf kwam.
\end{solution}

\begin{solution}{exo:g7:earthquakes:11}
(a) $400 \div 800 = 0.5$~uur: 30~minuten. (b) $20 \div 40 = 0.5$~uur:
30~minuten. (c) 60~minuten. (d) De laatste \qty{20}{km} duren even lang
als de eerste \qty{400}{km}: de golf vertraagt enorm in ondiep water.
(e) Het front van de golf vertraagt terwijl het water erachter nog snel
aankomt: het water stapelt zich op, en de golf wordt hoger.
\end{solution}

\begin{solution}{exo:g7:earthquakes:12}
(a) Drie stappen: $32 \times 32 \times 32 = 32\,768$, ongeveer
\textbf{33\,000}~keer. (b) $32\,768 \div 127 \approx 258$: ongeveer
260~jaar van alle \omterm{def:g4:volcanoes-earthquakes-first-look:earthquake}{aardbevingen} met \omterm{def:g7:earthquakes:magnitude}{magnitude} 6 ter wereld. (c) Bijna alle
energie komt vrij bij de paar allergrootste \omterm{def:g4:volcanoes-earthquakes-first-look:earthquake}{aardbevingen}: één enkele
\omterm{def:g4:volcanoes-earthquakes-first-look:earthquake}{aardbeving} met \omterm{def:g7:earthquakes:magnitude}{magnitude}~9 weegt op tegen eeuwen van \omterm{def:g4:volcanoes-earthquakes-first-look:earthquake}{aardbevingen} met
\omterm{def:g7:earthquakes:magnitude}{magnitude} 6.
\end{solution}

\begin{solution}{pb:g7:earthquakes:1}
\textbf{1.} A: $45 - 25 = 20$~s. B: $29 - 16 = 13$~s. C: $34 - 19 =
15$~s.

\textbf{2.} Station B: zijn vertraging is de kortste, en hoe korter de
vertraging, hoe dichterbij de \omterm{def:g4:volcanoes-earthquakes-first-look:earthquake}{aardbeving}. (Zijn snellere golven kwamen ook
als eerste aan.)

\textbf{3.} A: $8 \times 20 = \qty{160}{km}$. B: $8 \times 13 =
\qty{104}{km}$. C: $8 \times 15 = \qty{120}{km}$.

\textbf{4.} Alle punten op \qty{160}{km} van A: een cirkel met straal
\qty{160}{km} en middelpunt A. Het \omterm{def:g7:earthquakes:focus}{epicentrum} ligt ergens op die cirkel.

\textbf{5.} A: $160 \div 20 = 8$~hokjes. B: $104 \div 20 = 5.2$~hokjes.
C: $120 \div 20 = 6$~hokjes.

\textbf{6.} E ligt $3 + 6 = 9$~hokjes ten oosten en 14.5~hokjes ten
noorden van O, op \omterm{def:g2:water-on-earth:ocean}{zee} tussen de drie stations; de cirkels van A en B gaan
erdoorheen.

\textbf{7.} Onder de \omterm{def:g2:water-on-earth:ocean}{zee}. Een plotselinge beweging van de zeebodem kan een
\omterm{def:g7:earthquakes:tsunami}{tsunami} op gang brengen: de \omterm{def:g1:landscapes-around-us:water}{kusten} en het eiland moeten worden
gewaarschuwd.

\textbf{8.} Eén stap in \omterm{def:g7:earthquakes:magnitude}{magnitude}: ongeveer 32~keer meer energie.

\textbf{9.} $12 \times 20 = \qty{240}{km}$.

\textbf{10.} $240 \div 800 = 0.3$~uur, en $0.3 \times 60 = 18$~minuten.

\textbf{11.} 5~minuten na de \omterm{def:g4:volcanoes-earthquakes-first-look:earthquake}{aardbeving}.

\textbf{12.} $18 - 5 = 13$: de inwoners van Port Hale hebben ongeveer
\textbf{13~minuten} om hoger gelegen grond te bereiken. (De golf vertraagt
in het ondiepe water bij het eiland, wat iets meer tijd geeft, maar
niemand mag daarop rekenen: ze vertrekken meteen.)
\end{solution}
